% strong/confinement_topology.tex
% Part of the Hopfion Framework Repository
% Source: Paper XVIII — Preimage Topology and Confinement:
%   The Q_H=3 Sector of the Density-Feedback Hopfion
% Status: draft
% Last updated: 2026-07-04
%
%% hopfion-framework/preamble.tex
%% Shared preamble for all repository files.
%% \input this at the top of any standalone document.
%% Repository files are designed to be \input into papers directly.

\usepackage{amsmath,amssymb,amsthm}
\usepackage{hyperref}
\usepackage{enumitem}
\usepackage{booktabs}
\usepackage{xcolor}
\usepackage{geometry}
\geometry{margin=2.5cm}

%── Theorem environments ──────────────────────────────────────────────────────
\newtheorem{theorem}{Theorem}[section]
\newtheorem{lemma}[theorem]{Lemma}
\newtheorem{proposition}[theorem]{Proposition}
\newtheorem{corollary}[theorem]{Corollary}
\theoremstyle{definition}
\newtheorem{definition}[theorem]{Definition}
\newtheorem{result}[theorem]{Result}
\theoremstyle{remark}
\newtheorem{remark}[theorem]{Remark}
\newtheorem{observation}[theorem]{Observation}
\newtheorem{conjecture}{Conjecture}
\newtheorem{construction}[theorem]{Construction}
\newtheorem{condition}[theorem]{Condition}

%── Repository metadata commands ─────────────────────────────────────────────
% Usage: \status{published} or \status{open} or \status{conjectured}
\newcommand{\status}[1]{\par\noindent{\color{blue}\textbf{Status:} #1}}
% Usage: \source{Paper I, Theorem thm:scale}
\newcommand{\source}[1]{\par\noindent{\color{darkgray}\textbf{Source:} #1}}
% Usage: \residual{0.013\%} or \residual{exact}
\newcommand{\residual}[1]{\par\noindent{\color{teal}\textbf{Residual:} #1}}
% Usage: \depends{foundations/pentagon, foundations/suppression}
\newcommand{\depends}[1]{\par\noindent{\color{purple}\textbf{Depends on:} #1}}

%── Framework macros ─────────────────────────────────────────────────────────
\newcommand{\vp}{\varphi}
\newcommand{\bst}{\beta^{*}}
\newcommand{\Seff}{S_{\mathrm{eff}}}
\newcommand{\TCMB}{T_{\mathrm{CMB}}}
\newcommand{\Lcond}{\Lambda_{\mathrm{cond}}}
\newcommand{\LUV}{\Lambda_{\mathrm{UV}}}
\newcommand{\LQCD}{\Lambda_{\mathrm{QCD}}}
\newcommand{\MPl}{M_{\mathrm{Pl}}}
\newcommand{\ERy}{E_{\mathrm{Ry}}}
\newcommand{\aem}{\alpha_{\mathrm{em}}}
\newcommand{\sinW}{\sin^{2}\!\theta_{W}}
\newcommand{\vEW}{v_{\mathrm{EW}}}
\newcommand{\twoI}{2I}
\newcommand{\twoT}{2T}
\newcommand{\twoO}{2O}
\newcommand{\Ehat}{\hat{E}_8}
\newcommand{\QH}{Q_H}
\newcommand{\Ztwo}{\mathbb{Z}_2}
\newcommand{\Zthree}{\mathbb{Z}_3}
\newcommand{\SU}{\mathrm{SU}}
\newcommand{\rCMB}{\rho_{\mathrm{CMB}}}
\newcommand{\Lobs}{\Lambda_{\mathrm{obs}}}
\newcommand{\rinfty}{\rho_{\infty}}
\newcommand{\oBD}{\omega_{\mathrm{BD}}}
\newcommand{\xNMC}{\xi_{\mathrm{NMC}}}
\newcommand{\Efb}{E_{\mathrm{fb}}}
\newcommand{\Kfb}{K_{\mathrm{fb}}}
\newcommand{\dimq}[1]{d_{#1}}
\newcommand{\deff}{d_{\mathrm{eff}}}
\newcommand{\Kth}{K_{\mathrm{th}}}
\newcommand{\Ja}{\mathcal{J}_a}
\newcommand{\Jfour}{\mathcal{J}_4}
\newcommand{\Jiso}{\mathcal{J}_{\mathrm{iso}}}
\newcommand{\rhoinf}{\rho_{\infty}}
\newcommand{\rhoa}{\rho_{\mathrm{atom}}}
\newcommand{\Heff}{H_{\mathrm{eff}}}
\newcommand{\epsr}{\varepsilon(r)}
\newcommand{\Hcond}{H_{\mathrm{cond}}}
\newcommand{\eps}{\varepsilon}
\newcommand{\Rzero}{R_0}
\newcommand{\Cstar}{C^*}
\newcommand{\mH}{m_H}
\newcommand{\aem}{\alpha_{\mathrm{em}}}
\newcommand{\ms}{\mu^*}
\newcommand{\lam}{\lambda}
\newcommand{\fzero}{f_0}
\newcommand{\ftwo}{f_2}
\newcommand{\Itor}[1]{I_{#1}^{\mathrm{tor}}}
\newcommand{\Jfb}{\mathcal{J}_{\mathrm{fb}}}
\newcommand{\bz}{\beta_0}
\newcommand{\sopt}{s_{\mathrm{opt}}}
\newcommand{\leff}{\lambda_{\mathrm{eff}}}
\newcommand{\ord}{\operatorname{ord}}
\newcommand{\G}{\mathcal{G}}
\newcommand{\Ia}{I_{2a}}
\newcommand{\Ib}{I_{4}}
\newcommand{\kk}{\kappa}
\newcommand{\aInv}{\alpha^{-1}}
\newcommand{\mWZW}{m_{\mathrm{WZW}}}
\newcommand{\tang}{\dot\Gamma}
\newcommand{\Jones}{V}
\newcommand{\lk}{\mathrm{lk}}
\newcommand{\phistar}{\phi^*}
\newcommand{\CC}{\mathbb{C}}
\newcommand{\HH}{\mathbb{H}}
\newcommand{\OO}{\mathbb{O}}
\newcommand{\ZZ}{\mathbb{Z}}
\newcommand{\pcond}{p_{\mathrm{cond}}}
\newcommand{\pQM}{p_{\mathrm{QM}}}
\newcommand{\xcond}{x_{\mathrm{cond}}}
\newcommand{\CHSH}{\mathrm{CHSH}}
\newcommand{\Tsirelson}{2\sqrt{2}}
\newcommand{\RR}{\mathbb{R}}
\newcommand{\Hil}{\mathcal{H}}
\newcommand{\CQ}{\mathcal{C}_Q}
\newcommand{\CQr}{\mathcal{C}_Q^{\mathrm{red}}}
\newcommand{\bOmega}{\boldsymbol{\Omega}}
\newcommand{\dmu}{d\mu}
\newcommand{\nlvl}[1]{n\!\left(#1\right)}
\newcommand{\IE}{\mathrm{IE}}
\newcommand{\bsF}{{}_2F_1}




\section*{Preimage Topology, Jones-Polynomial Confinement, and Jet Structure of the $\QH=3$ Sector}
\label{sec:confinement_topology}

% Results extracted from Paper XVIII (draft)
% Scan date: 2026-07-04
\depends{strong/topology_qh3.tex; strong/colour.tex; quantum/neutrino.tex}

%════════════════════════════════════════════════════════════════════
% The Jones Polynomial at the WZW Root of Unity
%════════════════════════════════════════════════════════════════════

\begin{theorem}[Satellite Jones invariant of the trefoil]
\label{P18:thm:trefoil_satellite_jones}
Let $q_5=e^{2\pi i/5}$. The trefoil $T(2,3)$ is not a free-standing
knot in this framework: by Proposition~\ref{P18:prop:satellite} it is
the $(2,3)$-cable of one component of the $\QH=2$ Hopf link, with the
other (companion) component still present and linked. Computing the
Jones invariant of this configuration correctly --- via the
Kauffman-bracket/Temperley--Lieb cabling formula applied to the
diagram with the companion included --- gives
\begin{equation}
  \Jones_{\mathrm{sat}}(q_5) = \frac{q_5^3}{\vp}, \qquad
  |\Jones_{\mathrm{sat}}(q_5)| = \frac{1}{\vp}.
\end{equation}
This is irrational: the golden ratio is present in the trefoil's
satellite invariant at $q_5$, not absent from it.
\end{theorem}
\status{draft}
\source{P18:thm:trefoil_satellite_jones}
\residual{exact}

\begin{remark}[The free-standing evaluation, for comparison]
\label{P18:rem:trefoil_satellite_correction}
Evaluated as an abstract, free-standing knot (ignoring the satellite
embedding), the trefoil's Jones polynomial at $q_5$ is
$-q_5^{-1}+q_5^3+q_5$, of modulus $\sqrt{4-\vp}$ --- a different
irrational value from the satellite invariant above. Neither is an
integer, and the golden ratio is present in both.
\end{remark}
\status{draft}
\source{P18:rem:trefoil_satellite_correction}

\begin{proposition}[The trefoil within the satellite twist family]
\label{P18:prop:satellite_twist_family}
Varying the number of half-twists $n$ between the two cable strands
(companion linked once to each cable strand throughout; $n=3$ is the
trefoil), the modulus $|\Jones_{\mathrm{sat}}^{(n)}(q_5)|$ is periodic
in $n$ with period $10$ and satisfies the mirror symmetry
$|\Jones_{\mathrm{sat}}^{(n)}(q_5)|=|\Jones_{\mathrm{sat}}^{(10-n)}(q_5)|$.
At the four smallest cases:
\begin{align}
  n=3\ (\text{trefoil},\ \QH=3,\ \text{fused}): &\quad
    |\Jones_{\mathrm{sat}}(q_5)| = 1/\vp,\\
  n=4\ (\text{unfused, three components}): &\quad
    \Jones_{\mathrm{sat}}(q_5) = 1,\\
  n=5\ (\text{cinquefoil},\ \text{fused}): &\quad
    |\Jones_{\mathrm{sat}}(q_5)| = \vp,\\
  n=6\ (\text{unfused, three components}): &\quad
    \Jones_{\mathrm{sat}}(q_5) = -1.
\end{align}
Every fused configuration (odd $n$, a genuine torus knot) gives an
irrational, golden-ratio-related modulus; every unfused configuration
(even $n$, separate un-merged cable strands) gives a rational value.
\end{proposition}
\status{draft}
\source{P18:prop:satellite_twist_family}
\residual{exact}

\begin{remark}[The confinement criterion: fusion, not sector, and a phase relation between $\QH=2$ and $\QH=3$]
\label{P18:rem:richer_criterion}
The satellite twist family of Proposition~\ref{P18:prop:satellite_twist_family},
together with a direct computation of the (uncabled) Hopf link's own
Jones invariant, gives a confinement criterion that is different
from --- and more subtle than --- the original integer/irrational
framing.  We record three pieces of structure and one open question.

\textbf{(A) Confinement marks fusion, not sector.}
The rationality-class split tracks whether the cable strands
have \emph{fused} into a genuine torus knot, not which topological
sector is involved:
\begin{center}
\renewcommand{\arraystretch}{1.3}
\begin{tabular}{ccccl}
\toprule
$n$ & Components & $|\Jones_{\rm sat}(q_5)|$ & Rational? & Notes \\
\midrule
3 & 2 (fused, $T(2,3)$) & $1/\vp$ & no & $\QH=3$ trefoil \\
4 & 3 (unfused)          & $1$ & \textbf{yes} & \\
5 & 2 (fused, $T(2,5)$) & $\vp$ & no & cinquefoil \\
6 & 3 (unfused)          & $1$ & \textbf{yes} & [$V=-1$, so rational] \\
\bottomrule
\end{tabular}
\end{center}
Every fused (odd $n$) state computed is irrational; every unfused
(even $n$) state computed is rational.  Whether the physical
formation path from $\QH=2$ to $\QH=3$ actually passes through the
$n=4$ configuration, and whether its rationality is what produces the
flat-bottomed energy landscape of Paper~XVI~\cite{Paper16}
(Remark~\ref{P18:rem:saddle_difficulty}), is an open conjecture, not
established: nothing in this section identifies the twist parameter
$n$ with a physically continuous coordinate of the actual
condensate dynamics.

\textbf{(B) The Hopf link and the trefoil satellite carry the same
modulus, related by a phase.}  Direct computation of the (uncabled)
Hopf link's own Jones polynomial (not the disjoint-union factor, a
different quantity discussed below) gives
\begin{equation}
  |\Jones_{T(2,2)}(q_5)| \;=\; \frac{1}{\vp}
  \;=\; |\Jones_{\rm sat}(T(2,3))(q_5)|,
  \qquad
  \Jones_{\rm sat}(q_5) \;=\; q_5^{-1}\cdot\Jones_{T(2,2)}(q_5).
  \label{P18:eq:phase_relation}
\end{equation}
The two sectors carry the \emph{same} modulus; they are distinguished
by a single unit of the phase $q_5$, not by an inversion of the
golden ratio.  Whether this phase relation has independent physical
content, or how (if at all) it connects to the simple-current
condition $\dimq{3/2}=1$ (an established, independent confinement
criterion from Paper~XVII~\cite{Paper17}, Proposition~7.2), is not
established here.
(The disjoint-union factor $-(q_5^{1/2}+q_5^{-1/2})=-\vp$, used
elsewhere in this paper for multi-component reactions,
Theorem~\ref{P18:thm:jones_rule}, is a different quantity from
$\Jones_{T(2,2)}(q_5)$ and should not be conflated with it.)

\textbf{(C) The tower structure is $\{1/\vp, 1, \vp, -1, \ldots\}$.}
The fused states (odd $n$) have moduli alternating between $1/\vp$
and $\vp$ under the mirror symmetry
$|\Jones_{\rm sat}^{(n)}|=|\Jones_{\rm sat}^{(10-n)}|$;
the unfused states (even $n$) alternate between $V=1$ and $V=-1$.
This period is $10$; whether this is the same $10$ as the
quantum-group order $Q_{\rm gr}=2(k+2)=10$ of the framework
(Paper~I~\cite{Paper1}), or a coincidence of this specific tangle's
eigenvalue orders, is an open question and is not established.

\textbf{(D) Where $-1$ lives: the unfused sector at $n=6$.}
The value $V=-1$ that was originally claimed for the trefoil ($n=3$)
does occur in the family, but at $n=6$ (an unfused, three-component
configuration, the next rational step beyond the cinquefoil).
The classical identity $V_{T(2,3)}(i)=-1$ holds at the 4th root of
unity $t=i=e^{2\pi i/4}$ (the $k=2$ WZW level, the Arf invariant
identity $V_K(i)=(-1)^{{\rm Arf}(K)}$), which is unrelated to
the $k=3$ WZW structure used throughout this framework.
The precise placement of $-1$ at $n=6$ (unfused, rational) rather
than $n=3$ (fused, irrational) is itself consistent with the
fused/unfused rule of point~(A).
\end{remark}
\status{draft}
\source{P18:rem:richer_criterion}

\begin{remark}[Where the integer value $-1$ actually occurs]
\label{P18:rem:minus_one_at_n6}
The value $\Jones(q_5)=-1$ once (incorrectly) claimed for the trefoil
itself does occur in this family --- exactly, at $n=6$, an unfused
three-component configuration --- rather than at $n=3$.
\end{remark}
\status{draft}
\source{P18:rem:minus_one_at_n6}

%════════════════════════════════════════════════════════════════════
% Hopf Preimage Structure and the Torus Knot Tower
%════════════════════════════════════════════════════════════════════

\begin{definition}[Torus knot/link tower]
\label{P18:def:torus_tower}
The torus knot/link $T(2,n)$ is the closure of the two-strand braid
$\sigma_1^n\in B_2$, embedded on the standard torus
$\mathbb{T}^2\hookrightarrow S^3$, winding $2$ times around the
longitude and $n$ times around the meridian. For $n$ odd, $T(2,n)$
is a knot (one component); for $n$ even, a link (two components,
linking number $n/2$).
\end{definition}
\status{draft}
\source{P18:def:torus_tower}

\begin{proposition}[Preimage topology of the $T(2,n)$ tower]
\label{P18:prop:preimage_tower}
In the density-feedback FN framework, the minimal-action condensate
configuration in sector $\QH=N$ has preimage topology:
$\QH=1$: single unknotted circle $T(2,1)$;
$\QH=2$: two-component Hopf link $T(2,2)$, $\lk(\gamma_1,\gamma_2)=1$,
$\QH=2\cdot\lk=2$;
$\QH=3$: trefoil knot $T(2,3)$ (single knotted circle), whose three
self-crossings correspond structurally to the three quark crossing
vertices of Papers~XV--XVI.
\end{proposition}
\status{draft}
\source{P18:prop:preimage_tower}
\depends{strong/topology_qh3.tex}

\begin{remark}[Linking number convention and the $\QH$ formula]
\label{P18:rem:link_convention}
Equation $\QH=\sum_{a<b}\lk(\gamma_a,\gamma_b)$ holds in the
convention where the Hopf link $T(2,2)$ has
$\lk(\gamma_1,\gamma_2)=1$ and contributes $\QH=2$. The key
topological fact: $\QH=2$ preimage is a two-component Hopf link,
$\QH=3$ preimage is a single knotted circle (trefoil).
\end{remark}
\status{draft}
\source{P18:rem:link_convention}

\begin{definition}[Satellite knot]
\label{P18:def:satellite}
Let $K$ be a knot in $S^3$ and $P$ a knot in the solid torus
$D^2\times S^1$. The \emph{satellite knot} $S(K,P)$ is obtained by
embedding $D^2\times S^1$ tubularly around $K$ and taking the image
of $P$ under this embedding.
\end{definition}
\status{draft}
\source{P18:def:satellite}

\begin{proposition}[Trefoil as satellite of the Hopf link]
\label{P18:prop:satellite}
The trefoil $T(2,3)$ is the satellite of the unknot $U$ with pattern
the $(2,3)$-cable. Threading the second component $\gamma_2$ (the
$\QH=1$ unknot) through the solid-torus neighbourhood $N(\gamma_1)$
of the first Hopf-link component with the $(2,3)$-cable framing
($2$ times around the longitude, $3$ times around the meridian)
converts $T(2,2)\cup\gamma_2$ into the single-component trefoil
$T(2,3)$.
\end{proposition}
\status{draft}
\source{P18:prop:satellite}

\begin{remark}[Physical meaning of the satellite construction]
\label{P18:rem:satellite_physical}
The $\QH=3$ trefoil configuration is the unique topological
completion of the $\QH=2$ Hopf-link configuration obtained by
threading one additional circle (the $\QH=1$ unknot) through the
torus with the $(2,3)$-cable framing. This framing is the
minimal-action framing in the $\QH=3$ sector, dictated by the same
BPS saddle condition that fixes $\lambda_{\rm eff}^*=\vp^6/2^{1/3}$
(Paper~I).
\end{remark}
\status{draft}
\source{P18:rem:satellite_physical}
\depends{profile/bps.tex}

%════════════════════════════════════════════════════════════════════
% Topological Reaction Rules
%════════════════════════════════════════════════════════════════════

\begin{definition}[Rationality class]
\label{P18:def:rational_class}
A topological sector $L$ is in the \emph{rational class} at the WZW
root if $\Jones_L(q_5)\in\mathbb{Z}$; otherwise it is in the
\emph{irrational class}.
\end{definition}
\status{draft}
\source{P18:def:rational_class}

\begin{theorem}[Jones polynomial selection rule]
\label{P18:thm:jones_rule}
At the WZW root $q_5=e^{2\pi i/5}$, topological reactions in the
density-feedback condensate obey: (i)~charge conservation
$\sum_i Q_{H,i}$; (ii)~rationality factorisation --- if the initial
state is irrational-class and the final state contains a
rational-class component, at least one other final-state component
must be irrational.
\end{theorem}
\status{draft}
\source{P18:thm:jones_rule}
\residual{Proof status: (ii) is stated as a selection rule
whose complete proof requires a systematic multi-component Jones
polynomial framework (Paper18, Open Problem O1); the disjoint-union
normalisation is fixed by the Kauffman-bracket convention of
Paper~XVII.}

%════════════════════════════════════════════════════════════════════
% Forced Pion Emission
%════════════════════════════════════════════════════════════════════

\begin{theorem}[Pion emission is topologically forced]
\label{P18:thm:pion_forced}
The reaction $\QH=2+\QH=2\to\QH=3+\QH=1$ is the unique minimal
topological reaction that (i)~conserves Hopf charge ($2+2=3+1$);
(ii)~converts two irrational-class objects into one rational-class
object (the $\QH=3$ trefoil) with one irrational-class companion;
(iii)~minimises the number of final-state components. The $\QH=1$
product is identified with the pion ($\pi^\pm$ or $\pi^0$).
\end{theorem}
\status{draft}
\source{P18:thm:pion_forced}
\depends{quantum/neutrino.tex}

\begin{corollary}[Pion as topological spectator]
\label{P18:cor:pion_spectator}
In the reaction $\QH=2+\QH=2\to\QH=3+\QH=1$, the $\QH=1$ pion is the
preimage circle that does not participate in forming the trefoil.
Its internal dynamics are governed by $\SU(2)_1$ WZW theory with
Majorana statistics (Paper~XVII, Corollary~cor:majorana).
\end{corollary}
\status{draft}
\source{P18:cor:pion_spectator}
\depends{quantum/neutrino.tex}

\begin{remark}[$\QH=2+\QH=1$ reaction]
\label{P18:rem:qh2_qh1}
The complementary reaction $\QH=2+\QH=1\to\QH=3$ is the satellite
construction directly: the $\QH=1$ unknot threads through the
$\QH=2$ Hopf link with the $(2,3)$-cable framing, with no spectator
(all three preimage circles are consumed). Physically the
topological analogue of $e^-+\pi^+\to p$.
\end{remark}
\status{draft}
\source{P18:rem:qh2_qh1}

\begin{remark}[Higher-charge incoming objects and minimality of the listed channels]
\label{P18:rem:higher_charge_incoming}
Charge conservation permits any decomposition summing to $3$, but
the two listed channels ($2+2\to3+1$ and $2+1\to3$) are the
physically minimal cases: they involve only established stable
sectors ($\QH=1,2$), they are the lowest-energy threshold
reactions, and the $2+2\to3+1$ channel is the one identified with
real hadronic processes. No stable $\QH\geq4$ configuration is
established by the framework.
\end{remark}
\status{draft}
\source{P18:rem:higher_charge_incoming}

\begin{remark}[Kinematic path of the $\QH=2\to\QH=3$ transition]
\label{P18:rem:crossing_path}
The $\QH=2$ Hopfion is two linked tubes ($\gamma_1$ outer,
$\gamma_2$ inner), not one reconfiguring tube. The transition
proceeds in three steps: (1)~writhe injection deforms $\gamma_2$
without changing $\lk(\gamma_1,\gamma_2)$; (2)~$\gamma_2$ threads
into $N(\gamma_1)$, forced asymmetric by the $(2,3)$-cable framing
constraint; (3)~once threading commences, all three trefoil
crossings are generated simultaneously as continuous consequences
of the completed cable winding --- a single continuous geometric
process, not three discrete crossing choices. The crossing vertex
formed is simultaneously the site of maximum $S_{\rm eff}$ and
minimum $\rho_{J_4}$ (Proposition~\ref{P18:prop:anticorrelation}). A
candidate explanation --- not established, offered as a conjecture
--- is that the crossing has no stable irrational algebraic anchor
because the formation event passes through the unfused $n=4$
intermediate (rational satellite invariant
$\Jones_{\mathrm{sat}}^{(4)}(q_5)=1$,
Proposition~\ref{P18:prop:satellite_twist_family}) before fusing to
the irrational $n=3$ value ($|\Jones_{\mathrm{sat}}^{(3)}(q_5)|=1/\vp$,
Theorem~\ref{P18:thm:trefoil_satellite_jones}); whether the formation
path genuinely passes through $n=4$ is not established. What
\emph{is} established is that the flat energy plateau itself is an
observed numerical fact (Papers~XV/XVI) and that the trefoil's own
satellite invariant is irrational: on either reading, the crossing
is immediately trying to reverse, and the baryon is metastable in
the dynamical, not energetic, sense.
\end{remark}
\status{draft}
\source{P18:rem:crossing_path}
\depends{strong/topology_qh3.tex}

\begin{remark}[Two sub-thresholds and the virtual baryon regime]
\label{P18:rem:two_thresholds}
The energy threshold governing $\QH=2\to\QH=3$ formation has two
internal sub-thresholds: (I)~the \emph{commitment} threshold, the
minimum kinetic energy to drive the fold/crossing of
Remark~\ref{P18:rem:crossing_path} against the density-feedback
restoring force (below it the field elastically restores $\QH=2$);
(II)~the \emph{ejection} threshold, the minimum energy for the
implosion--explosion cascade (Remark~\ref{P18:rem:implode_explode}) to
run to completion before the committed crossing reverses. Between
the two lies a \emph{virtual baryon} regime: $\QH=3$ is momentarily
achieved but reverses without ejecting a $\QH=1$ daughter.
\end{remark}
\status{draft}
\source{P18:rem:two_thresholds}

\begin{observation}[Geometric precursor at $q=2.3$]
\label{P18:obs:precursor}
In the $T(2,q_{\rm pol})$ construction sweep at fixed $\QH=2$ field
winding ($\alpha_{\rm wind}=0$), a sharp field-orientation
discontinuity already appears at $q_{\rm pol}=2.3$ at a single site,
purely geometric (Whitehead $\QH$ remains $2$). This single site is
the geometric precursor of the eventual topological threading site:
geometry precedes topology.
\end{observation}
\status{draft}
\source{P18:obs:precursor}
\residual{numerical (single simulated sweep, $N=64$, $h=0.175$)}

\begin{remark}[Production conditions: a two-threshold structure]
\label{P18:rem:production_conditions}
Two independent physical conditions govern $\QH=3$ formation:
(i)~an \emph{energy threshold}
$\Delta E_{2+2}=E(\QH{=}3)+E(\QH{=}1)-2E(\QH{=}2)>0$ (or the
analogous $\Delta E_{2+1}$), both reactions being kinematically
closed at rest since the baryon mass scale exceeds the
charged-lepton and pion scales; (ii)~a \emph{rate/density threshold}
$\Gamma\propto n^2\langle\sigma v\rangle$ (or $n_{\QH=2}n_{\QH=1}$).
Ordinary vacuum satisfies neither; collider beams, the primordial
quark-gluon plasma, and compact-object extremes satisfy both.
\end{remark}
\status{draft}
\source{P18:rem:production_conditions}
\residual{Quantitative evaluation of $\Delta E_{2+2}$, $\Delta E_{2+1}$
from the Paper~XVII mass formulas is Paper18 Open Problem O6.}

%════════════════════════════════════════════════════════════════════
% The phi* Prediction from T(2,3) Geometry
%════════════════════════════════════════════════════════════════════

\begin{theorem}[Crossing geometry of $T_{2,3}$]
\label{P18:thm:crossing_geometry}
For the standard trefoil parametrisation
$\Gamma(t)=((R_0+r_0\cos3t)\cos2t,(R_0+r_0\cos3t)\sin2t,r_0\sin3t)$:
(i)~the tangent $\dot\Gamma(t_n)$ lies exactly in the $z=0$ plane at
each crossing $t_n=\pi/6+2\pi n/3$; (ii)~the azimuthal angles of
$\dot\Gamma(t_n)$ are separated by exactly $2\pi/3$ for $n=0,1,2$;
(iii)~each pair of adjacent crossings is separated by exactly
$2\pi/3$ in azimuth.
\end{theorem}
\status{draft}
\source{P18:thm:crossing_geometry}
\residual{exact}
\depends{strong/topology_qh3.tex}

\begin{theorem}[$\phistar$ distribution for baryon-initiated production]
\label{P18:thm:phi_star}
When a $\QH=3$ baryon fragments to produce a $\rho^+\rho^-$ pair at
two of its three crossing vertices, with decay-plane orientation set
by the tangent direction at the production crossing, the
acoplanarity-angle distribution is
\begin{equation}
  W(\phistar) = \tfrac{1}{2\pi}\bigl[1+a_3\cos(3\phistar)\bigr],
  \qquad a_1=a_2=0,
\end{equation}
with $a_3\neq0$, peaked at $\phistar=2\pi/3$ and its $\Zthree$
images.
\end{theorem}
\status{draft}
\source{P18:thm:phi_star}
\residual{prediction (untested)}

\begin{remark}[Contrast with photon-initiated production]
\label{P18:rem:photon_contrast}
For $e^+e^-\to\gamma^*\to\rho^+\rho^-$ (no $\QH=3$ parent), the
$\phistar$ distribution is governed by transverse photon helicity:
$W_\gamma(\phistar)=\tfrac{1}{2\pi}[1+a_2^\gamma\cos(2\phistar)]$,
$a_1=a_3=0$. The distinction between $\cos(3\phistar)$
(baryon-initiated) and $\cos(2\phistar)$ (photon-initiated) is
unambiguous and parameter-free.
\end{remark}
\status{draft}
\source{P18:rem:photon_contrast}

\begin{proposition}[Three-jet coplanarity from B-departure geometry]
\label{P18:prop:coplanarity}
The three primary jets from a $\QH=3$ baryon exit at the three
B-segment departure zones, immediately past the crossing vertices;
by Theorem~\ref{P18:thm:crossing_geometry}(i) their tangents lie exactly
in the $z=+r_0$ plane, separated by $2\pi/3$. The three primary jets
therefore lie in a common elevated plane, giving zero intrinsic
acoplanarity at the topological level.
\end{proposition}
\status{draft}
\source{P18:prop:coplanarity}
\depends{strong/topology_qh3.tex}

%════════════════════════════════════════════════════════════════════
% Connection to the Arc-Segment Energy Partition
%════════════════════════════════════════════════════════════════════

\begin{proposition}[A-segment dominance]
\label{P18:prop:arcseg_dominance}
The A-segments (midpoint-to-crossing, 90°) carry $\approx79\%$ of
the total $J_4$ charge at energy density $\rho_A/\rho_B\approx1.23$
(Paper~XVI, Table~1, $N=96$ grid). The Hopf flux concentrates in the
propagation arms between midpoints and crossings, not at the
crossing vertices.
\end{proposition}
\status{draft}
\source{P18:prop:arcseg_dominance}
\depends{strong/topology_qh3.tex}

\begin{proposition}[$\rho_{J_4}$ anti-correlates with $S_{\mathrm{eff}}$]
\label{P18:prop:anticorrelation}
$S_{\rm eff}(\theta,\rho)=\sin^4\theta/[\vp^6(1+\beta\rho)]$ is
maximised at the crossing vertices ($\kappa=0.399$) and minimised at
the midpoints ($\kappa=0.026$). The field $\rho_{J_4}$ is depleted
precisely where $S_{\rm eff}$ is largest: crossing vertices are
topological bottlenecks. This anti-correlation is the
condensate-level signature of confinement.
\end{proposition}
\status{draft}
\source{P18:prop:anticorrelation}
\depends{strong/topology_qh3.tex}

\begin{proposition}[Energy cascade through the $T_{2,3}$ tube]
\label{P18:prop:cascade_p18}
The physical energy cascade in the $\QH=3$ condensate ($N=256$
gradient-flow run, Paper~XVI) follows: (i)~initial fraying at the
distal lobes (maximum radial exposure); (ii)~crossing depletion
propagates inward (the crossing cannot radiate outward since its
tangent lies exactly horizontal; it acts as a topological lens);
(iii)~redirection at B-segment departure produces the three primary
jets, coplanar at $z=+r_0$; (iv)~secondary lobe emission (jets 4--6)
deflected $\approx16$--$19^\circ$ from the three-jet plane;
(v)~midpoints show only minor perturbation throughout, remaining
most protected. The crossing vertex is an energy redirector, not an
energy source.
\end{proposition}
\status{draft}
\source{P18:prop:cascade_p18}
\depends{strong/topology_qh3.tex}

\begin{remark}[Implosion, then explosion]
\label{P18:rem:implode_explode}
The $N=256$ snapshots support a two-stage decay picture: the three
lobes deplete together diffusely and the displaced energy is
reflected \emph{inward} along the A-segment (implosion), because the
topological floor forbids direct radial escape; later the field
reorganises into sharply bounded bright arcs as the inward energy is
redirected through the crossings into the B-departure channels
(explosion). The crossings are conduits throughout, never sources.
The Kleckner--Irvine trefoil vortex knots in water offer only a
partial analogy: reconnection there is triggered by strand proximity
at the crossings, whereas the condensate's depletion begins at the
lobes (radial exposure to vacuum), a genuine point of disanalogy
traceable to the density-feedback term $\sin^4(f)$.
\end{remark}
\status{draft}
\source{P18:rem:implode_explode}

\begin{remark}[Cascade ordering]
\label{P18:rem:cascade_ordering}
The operative weakness driving the cascade order is \emph{radial
exposure} (distance from the tube axis relative to the vacuum
boundary), not curvature: midpoints ($R_{xy}=R_0-r_0$) are
geometrically unconstrained but radially protected; distal lobes
($R_{xy}=R_0+r_0$) are geometrically intermediate but radially
exposed and shed energy first under gradient flow.
\end{remark}
\status{draft}
\source{P18:rem:cascade_ordering}

\begin{remark}[Why the saddle point is elusive in Paper~XVI]
\label{P18:rem:saddle_difficulty}
At $\QH=2$, the BPS condition $\Kfb/J_4=\vp^6/2^{1/3}$ pins a unique
irrational profile fixed point. At $\QH=3$, the trefoil's own
satellite Jones invariant is irrational
($|\Jones_{\mathrm{sat}}(q_5)|=1/\vp$,
Theorem~\ref{P18:thm:trefoil_satellite_jones}), so this argument does
not, on its own, explain the elusive saddle. One conjectural (not
established) candidate explanation is that the physical formation
path passes through the unfused $n=4$ intermediate, which has a
rational satellite invariant
($\Jones_{\mathrm{sat}}^{(4)}(q_5)=1$,
Proposition~\ref{P18:prop:satellite_twist_family}); if so, it would
be the rationality of this intermediate, not the irrational value of
the completed trefoil, that removes the irrational algebraic anchor
during formation. The energy landscape's flat plateau is itself an
established numerical fact (Papers~XV/XVI gradient-flow searches);
gradient flow converges to this family (explaining the reproducible
arc-segment partition and $\rho_{J_4}$--$S_{\rm eff}$
anti-correlation) without converging to a specific point within it
(explaining why a unique $C^*$ is not found). Consequently the minor
radius $r_0=0.874$ is set by the profile-construction ansatz, not a
BPS minimisation; two candidate exact forms,
$r_0=2R_0/(3\vp+2)=0.8754$ (Paper~XVII solar-angle match) and
$r_0=R_0/C_2^*=0.8742$ (satellite-construction match), are shown
self-consistent with $\lambda_3=\vp^6$ via
$\lambda_3\cdot r_0=2R_0\vp^2$, but the flat landscape cannot decide
between them.
\end{remark}
\status{draft}
\source{P18:rem:saddle_difficulty}
\depends{strong/topology_qh3.tex; quantum/neutrino.tex}

\begin{remark}[Status of the cascade-ordering interpretation]
\label{P18:rem:three_remark_interpretation}
$E_{\rm geom}=K\cdot J_4$ is a global sum; the redistribution
described by Proposition~\ref{P18:prop:cascade_p18} is genuine global
dynamics, not a per-segment artefact. The revised ordering does not
invalidate the six-quark conjecture (Conjecture~\ref{P18:conj:six_quarks}):
midpoints remain the geometrically weakest sites under Frenet
analysis and the expected breaking sites under thermal/quantum
fluctuations; the gradient-flow ordering instead reflects the
steepest-descent path from the family boundary. The near-perfect
$\Zthree$ symmetry of the $N=256$ run confirms that the $(2{+}1)$
breaking seen at $N=192$ was a numerical (Adam-optimiser) artefact.
\end{remark}
\status{draft}
\source{P18:rem:three_remark_interpretation}

\begin{remark}[Two-phase landscape, $\Zthree$-symmetry breaking, and the Jones-polynomial explanation]
\label{P18:rem:gf_twophase_p18}
The $N=216$, $h=0.05$ run refines the sector-boundary estimate to
$K_{\min}\approx2289$ (from $2294$ at $N=192$). Phase~1 genuinely
operates within the $\QH=3$ sector (field sharpens, no topology
change). Phase~2 exposes a structural, non-numerical feature: the
field spontaneously breaks the exact $\Zthree$ symmetry of the
per-strand construction because it operates on the flat-bottomed
valley of Remark~\ref{P18:rem:saddle_difficulty} (an established
numerical fact; the conjectural rational-$n=4$-intermediate account
of why the valley is flat, discussed there, remains unestablished).
Floating-point accumulation in the Adam optimiser breaks $\Zthree$
at the $\sim10^{-7}$ level per step; the resulting $(2{+}1)$
splitting matches the single-midpoint string-breaking mode of
Conjecture~\ref{P18:conj:six_quarks}.
\end{remark}
\status{draft}
\source{P18:rem:gf_twophase_p18}
\depends{strong/topology_qh3.tex}

\begin{remark}[$K_{\min}$ resolution convergence and three-point extrapolation]
\label{P18:rem:kmin_extrapolation}
$K_{\min}$ at $N=144,192,256$ ($h=0.075,0.05,0.0375$) is
$2313.9$, $2289$, $2268$. Fitting
$K_{\min}(h)\approx K_{\min}^\infty+Ah^\alpha$ gives
$\alpha\approx0.75\pm0.15$ and
$K_{\min}^\infty\approx2230$--$2245$, monotone convergence from
above consistent with $O(h^\alpha)$ discretisation error.
\end{remark}
\status{draft}
\source{P18:rem:kmin_extrapolation}
\residual{numerical, three-point extrapolation}

%════════════════════════════════════════════════════════════════════
% The Botner Measurement Framework
%════════════════════════════════════════════════════════════════════

\begin{proposition}[Topological $\phistar$ distribution in Botner form]
\label{P18:prop:botner_analog}
By analogy with the DELPHI Botner three-jet form
$(1+\beta\cos2\chi)/2\pi$, the $\phistar$ distribution for
baryon-initiated $\rho^+\rho^-$ production is
$(1+\beta_3\cos3\phistar)/2\pi$, with
\begin{equation}
  \beta_3^{\rm est}\sim\frac{J_4^A-J_4^B}{J_4^A+J_4^B}
  =\frac{79\%-21\%}{79\%+21\%}=0.58,
\end{equation}
using $J_4^A/J_4^B\approx3.76$ from Paper~XVI. This is a
zero-free-parameter geometric prediction: photon/Z-initiated
production gives $\cos2\chi$ ($\Ztwo$, spin-1 photon);
baryon-initiated gives $\cos3\phistar$ ($\Zthree$, $T_{2,3}$
trefoil).
\end{proposition}
\status{draft}
\source{P18:prop:botner_analog}
\residual{prediction, $\beta_3\approx0.58$ (untested)}
\depends{strong/topology_qh3.tex}

%════════════════════════════════════════════════════════════════════
% Four-Jet Out-of-Plane Structure and the Jeremie Anomaly
%════════════════════════════════════════════════════════════════════

\begin{remark}[Two distinct orderings: depletion order versus jet hardness]
\label{P18:rem:two_orderings}
The condensate-internal depletion order (lobes first, crossings
last to show depletion, midpoints avoided) is a different question
from the collider jet-hardness order (B-departure exits hardest,
then lobes, then softer categories): the crossing vertex is a
redirector throughout (Proposition~\ref{P18:prop:cascade_p18}) and never
releases energy directly, but the hardest jets still exit at the
B-departure zones it redirects to.
\end{remark}
\status{draft}
\source{P18:rem:two_orderings}

\begin{proposition}[Out-of-plane deflection of the 4th jet]
\label{P18:prop:fourjet_oop}
The 12-site fraying pattern divides into four $\Zthree$-symmetric
groups of three. Any single group produces three symmetric jets
separated by $2\pi/3$; the 4th jet is necessarily the first emission
from a geometrically distinct second group. Under the ordering of
Proposition~\ref{P18:prop:jet_bound} (jets 1--3 at B-departure exits,
jet~4 at the first distal lobe, tangent $z$-component
$+0.321|\dot\Gamma|$), the deflection is
$\Delta\theta_{\rm lobe}\approx16^\circ$--$19^\circ$ below the
three-jet plane. A broad, unimodal QCD-like distribution is
inconsistent with any alternating-group ordering.
\end{proposition}
\status{draft}
\source{P18:prop:fourjet_oop}
\residual{prediction, $16^\circ$--$19^\circ$ (untested)}
\depends{strong/topology_qh3.tex}

\begin{remark}[The Jeremie anomaly]
\label{P18:rem:jeremie}
Jeremie et al.\ (OPAL, LEP) found shape differences between
four-jet Bengtsson--Zerwas angle data and ERT QCD matrix-element
predictions concentrated at $m_3+m_4>10\,\mathrm{GeV}$. This is
consistent with the framework's prediction
(Proposition~\ref{P18:prop:fourjet_oop}) of a discrete, topologically
distinct 4th jet at $16^\circ$--$19^\circ$ out of plane, with
tube-carried non-perturbative jets expected to carry higher
invariant mass than perturbative collinear gluons.
\end{remark}
\status{draft}
\source{P18:rem:jeremie}
\residual{consistent with existing LEP/OPAL data; not a dedicated test}

\begin{proposition}[Jet multiplicity upper bound and energy ordering]
\label{P18:prop:jet_bound}
The 12-site fraying pattern of Paper~XVI gives a topological upper
bound of 12 distinct ejection events per $\QH=3$ baryon, in four
groups of three: B-departure exits (jets 1--3, hardest, coplanar at
$z=+r_0$); distal-lobe secondaries (jets 4--6,
$16^\circ$--$19^\circ$ off-plane); crossing approach-side entry
disturbances (jets 7--9, softer); midpoint torsion-twist residuals
(jets 10--12, softest). The 12-jet bound and the midpoints-last
ranking are robust to the specific ordering of the middle two
groups.
\end{proposition}
\status{draft}
\source{P18:prop:jet_bound}
\residual{prediction: rate vanishes above 12 jets per baryon (untested,
requires $\gtrsim500\,\mathrm{GeV}$ lepton collider)}
\depends{strong/topology_qh3.tex}

%════════════════════════════════════════════════════════════════════
% Speculative Correspondences (flagged as conjectures)
%════════════════════════════════════════════════════════════════════

\begin{conjecture}[Six-quark combinatorial]
\label{P18:conj:six_quarks}
The $\QH=3$ baryon decays by losing topological integrity at one or
more of its three midpoint torsion-twist sites. The six quark
flavours correspond to the six distinct non-trivial breaking
patterns $\binom{3}{1}+\binom{3}{2}=3+3=6$: the three single-breaking
patterns and the three double-breaking patterns forming two isospin
triplets (charge $+2/3$ and $-1/3$, or vice versa).
\end{conjecture}
\status{draft}
\source{P18:conj:six_quarks}
\depends{strong/quark_masses.tex}

\begin{remark}[Constraints on any future derivation]
\label{P18:rem:quark_constraint}
Any derivation of the quark mass spectrum within this framework must
(i)~reproduce the $3+3$ combinatorial pattern; (ii)~assign the
correct isospin to each triplet; (iii)~derive the energy cost of
each breaking pattern from the arc-segment $J_4$ deficit $\delta J_4$
at each midpoint relative to the adjacent A-segment average.
\end{remark}
\status{draft}
\source{P18:rem:quark_constraint}
\depends{strong/quark_masses.tex}

\begin{conjecture}[Writhe asymmetry and isospin]
\label{P18:conj:isospin}
The trefoil $T_{2,3}$ has an exact $\Ztwo$ writhe asymmetry: the
three crossing vertices sit at $z=+r_0$, while the three midpoints
and three outer-lobe distal points sit at $z=0$. The isospin
assignment $I_3=+1/2$ (charge $+2/3$) corresponds to the
crossing-vertex network; $I_3=-1/2$ (charge $-1/3$) corresponds to
the midpoint/distal-lobe network.
\end{conjecture}
\status{draft}
\source{P18:conj:isospin}
\depends{strong/topology_qh3.tex}

\begin{proposition}[Isospin-triplet mass-scale ratio from tangent geometry]
\label{P18:prop:isospin_mass_ratio}
Defining the out-of-plane fraction $\dot\Gamma_z^2$ at each site
category (zero at crossings; $\langle\dot\Gamma_z^2\rangle_{M,L}
=\tfrac12(0.525^2+0.321^2)=0.189$ at midpoints/lobes), the
zero-free-parameter estimate
$M_{\rm down}/M_{\rm up}\sim0.189$ compares to the PDG geometric-mean
value $0.157$ (up: $u,c,t$; down: $d,s,b$), a $21\%$ discrepancy ---
comparable to, and better than, several of the mechanisms ruled out
in Paper~XVI's systematic survey.
\end{proposition}
\status{draft}
\source{P18:prop:isospin_mass_ratio}
\residual{21\% (untested first-principles derivation)}
\depends{strong/quark_masses.tex; strong/topology_qh3.tex}

\begin{remark}[Scope of the isospin mass-ratio estimate]
\label{P18:rem:isospin_mass_ratio_scope}
The identification of $M_{\rm down}/M_{\rm up}$ with the
out-of-plane/in-plane tangent-fraction ratio is motivated by the
writhe asymmetry of Conjecture~\ref{P18:conj:isospin} but not derived
from the density-feedback action. It does not address the
intra-triplet hierarchy ($m_u\ll m_c\ll m_t$, etc.), which requires
either the generation-dependent string-tension contribution
(Paper~XVI \S12) or a dynamical treatment of the
$\QH=2+\QH=1\to\QH=3$ formation event.
\end{remark}
\status{draft}
\source{P18:rem:isospin_mass_ratio_scope}

%════════════════════════════════════════════════════════════════════
% Dynamical Origin of Quark Masses — Paper XIX (draft)
%════════════════════════════════════════════════════════════════════

\begin{lemma}[Profile slope at the transition midpoint]
\label{P19:lem:profile_slope}
For the $\QH=2$ profile function $f(\rho)=2\arctan(\rho^{-C})$,
$f'(1)=-C$.
\end{lemma}
\status{draft}
\source{P19:lem:profile_slope}
\depends{strong/topology_qh3.tex}

\begin{theorem}[Analytic derivation of the trefoil minor radius]
\label{P19:thm:r0_derivation}
The trefoil minor radius equals the $\QH=2$ profile wall width:
$r_0 = R_0/C_2^*$. The satellite embedding
(Proposition~\ref{P18:prop:satellite}) requires the trefoil's minor
radius to equal the radial extent of the $\QH=2$ domain wall for the
embedding to be non-singular. Numerically $R_0/C_2^*=0.874172$,
matching the numerical construction value $r_0=0.874$ to $0.02\%$.
This closes the open problem, shared by Papers~XV and~XVIII, of the
physical origin of $r_0$.
\end{theorem}
\status{draft}
\source{P19:thm:r0_derivation}
\residual{0.02\%}
\depends{strong/confinement_topology.tex; strong/topology_qh3.tex}

\begin{remark}[Consistency with the solar-angle candidate of Paper~XVII]
\label{P19:rem:solar_angle_reconcile}
Paper~XVII's alternative candidate $r_0=2R_0/(3\vp+2)=0.8754$
(a $0.16\%$ match, motivated by the midpoint emission angle equalling
$\arctan(1/\vp)$) is a numerical coincidence of the near-equality
$C_2^*\approx\vp R_0/\sqrt2$ at $R_0=3$, not an independent derivation.
Theorem~\ref{P19:thm:r0_derivation}'s $r_0=R_0/C_2^*$ (a $0.02\%$
match, a factor of eight closer) is the derived value.
\end{remark}
\status{draft}
\source{P19:rem:solar_angle_reconcile}
\depends{quantum/neutrino.tex}

\begin{corollary}[The Bogomolny consistency $\lambda_3\cdot r_0=2R_0\vp^2$]
\label{P19:cor:lambda3_consistency}
Combining Theorem~\ref{P19:thm:r0_derivation} with the proved $\QH=3$
Bogomolny parameter $\lambda_3=\vp^6$, $\lambda_3\cdot r_0=\vp^6 R_0/C_2^*$
holds identically; numerically the right side ($15.68$) matches
$2R_0\vp^2=15.71$ to the same $0.16\%$ residual as the solar-angle
comparison of Remark~\ref{P19:rem:solar_angle_reconcile}, reflecting
the same underlying near-equality rather than an independent
constraint.
\end{corollary}
\status{draft}
\source{P19:cor:lambda3_consistency}
\residual{0.16\%}
\depends{strong/topology_qh3.tex}

%────────────────────────────────────────────────────────────────────
% Reconciling the six-quark and isospin conjectures
%────────────────────────────────────────────────────────────────────

\begin{remark}[Paper~XIX restates both Paper~XVIII conjectures before reconciling them]
\label{P19:conj:six_quarks_r19}
\label{P19:conj:isospin_r19}
Paper~XIX restates Conjecture~\ref{P18:conj:six_quarks} (six-quark
combinatorial) and Conjecture~\ref{P18:conj:isospin} (writhe-isospin)
verbatim under its own labels \texttt{conj:six\_quarks\_r19} and
\texttt{conj:isospin\_r19} before identifying that they are
internally inconsistent as separately stated: the former attributes
all six quark flavours to midpoint breaking, while the latter assigns
midpoints exclusively to the $I_3=-\tfrac12$ triplet, so only three
(not six) quarks could arise from midpoint breaking under the
isospin assignment. This inconsistency motivates the reformulation
below.
\end{remark}
\status{draft}
\source{P19:conj:six_quarks_r19,P19:conj:isospin_r19}

\begin{conjecture}[Quark flavours as partition channels]
\label{P19:conj:six_quarks_reformulated}
The six quark flavours correspond to the six characteristic scales of
the dynamical partition of the incoming perturbation energy
$E_{\mathrm{pert}}$ across the two networks of trefoil sites: three
\emph{primary channels} through the crossing/B-departure network
($z=+r_0$), giving the $I_3=+\tfrac12$ triplet $(u,c,t)$; and three
\emph{residual channels} through the midpoint/lobe network ($z=0$),
giving the $I_3=-\tfrac12$ triplet $(d,s,b)$. Within each triplet the
mass hierarchy comes from the azimuthal position of the site relative
to the threading site of the formation event; the two-triplet overall
scale ratio comes from the writhe-$\Ztwo$-weighted tangent-geometry
difference between the two networks. This reconciles
Conjecture~\ref{P18:conj:six_quarks} (six-quark combinatorial) and
Conjecture~\ref{P18:conj:isospin} (writhe-isospin), which are
internally inconsistent as separately stated: the former attributes
all six flavours to midpoint breaking, while the latter assigns
midpoints exclusively to $I_3=-\tfrac12$.
\end{conjecture}
\status{draft}
\source{P19:conj:six_quarks_reformulated}
\depends{strong/quark_masses.tex}

\begin{remark}[Consistency check with the observed sextet]
\label{P19:rem:sextet_check}
The reformulation retains the count $3+3=6$ but attributes it to
network membership rather than to a combinatorial choice of breaking
pattern; Conjecture~\ref{P18:conj:six_quarks}'s
$\binom31+\binom32=6$ counting is recovered as the combinatorial
signature of two independent $\Zthree$-broken triplets, and the
isospin assignment of Conjecture~\ref{P18:conj:isospin} is preserved.
\end{remark}
\status{draft}
\source{P19:rem:sextet_check}

%────────────────────────────────────────────────────────────────────
% Falsification of the single-loss-coefficient cascade model
%────────────────────────────────────────────────────────────────────

\begin{proposition}[Single-loss cascade is falsified]
\label{P19:prop:cascade_falsified}
Neither the arc-segment density ratio $\eta_B=\rho_B/\rho_A=0.81$ nor
the per-crossing transmission $\eta_{\mathrm{cross}}=0.266$ (both from
Paper~XVI) reproduces observed intra-triplet quark mass ratios under
the single-loss-coefficient cascade $E_1:E_2:E_3=1:\eta:\eta^2$: the
predicted second-position ratio misses observation by factors of
$37$--$109$, and the third-position ratio by five orders of
magnitude, for both isospin triplets.
\end{proposition}
\status{draft}
\source{P19:prop:cascade_falsified}
\depends{strong/topology_qh3.tex}

\begin{remark}[What the falsification does and does not rule out]
\label{P19:rem:cascade_scope}
Proposition~\ref{P19:prop:cascade_falsified} rules out the specific
model in which the entire intra-triplet mass hierarchy is generated by
a single per-crossing propagation loss coefficient; it does not rule
out cascade physics as such (Proposition~P18:prop:cascade\_p18), since
additional generation-dependent physics (e.g.\ $\QH=3$-internal
string-tension energy, or different midpoint-network propagation) may
still be required.
\end{remark}
\status{draft}
\source{P19:rem:cascade_scope}

%────────────────────────────────────────────────────────────────────
% Falsification of the uniform phi-tower generation multiplier
%────────────────────────────────────────────────────────────────────

\begin{definition}[Uniform $\vp$-tower generation multiplier]
\label{P19:def:phi_tower}
The hypothesis that the generation-dependent contribution to the
quark mass at generation $g\in\{1,2,3\}$ takes the form
$\Delta m^{(g)}=M_0\vp^{2(g-1)a}$ for some integer $a\ge1$, giving
intra-triplet ratios $\Delta m^{(2)}/\Delta m^{(1)}=
\Delta m^{(3)}/\Delta m^{(2)}=\vp^{2a}$.
\end{definition}
\status{draft}
\source{P19:def:phi_tower}

\begin{proposition}[Uniform $\vp$-tower multiplier is falsified]
\label{P19:prop:phi_tower_falsified}
No integer $a$ satisfies $m^{(2)}/m^{(1)}=m^{(3)}/m^{(2)}=\vp^{2a}$
(Definition~\ref{P19:def:phi_tower}) for either the up-type or the
down-type triplet: the consecutive $\vp$-exponent gaps are
$(13.26,10.20)$ for up-type and $(6.23,7.90)$ for down-type, and the
two triplets require different values of $a$ even under a
most-generous fit.
\end{proposition}
\status{draft}
\source{P19:prop:phi_tower_falsified}
\depends{strong/quark_masses.tex}

\begin{remark}[The observed exponent pattern shows internal structure]
\label{P19:rem:tower_structure}
The averages of the two exponent gaps per triplet are $11.73$
(up-type) and $7.06$ (down-type), whose ratio $1.661$ approximates
$\vp=1.618$ to $3\%$: the up-type tower is stretched relative to the
down-type tower by roughly a factor of $\vp$ in exponent.
\end{remark}
\status{draft}
\source{P19:rem:tower_structure}

\begin{proposition}[The Hopf-link/quantum-dimension ratio does not account for this ratio]
\label{P19:prop:jones_modulus_ratio}
The up-type network's coupling weight is the Hopf link's Jones
modulus $|\Jones_{T(2,2)}(q_5)|=1/\vp$
(Theorem~\ref{P19:thm:channel_1}); the down-type network's coupling
weight is the quantum dimension $\dimq{3/2}=1$. Their ratio
$|\Jones_{T(2,2)}(q_5)|/\dimq{3/2}=1/\vp\approx0.618$ is off from the
observed generation-spacing ratio $\bar n_{\mathrm{up}}/\bar
n_{\mathrm{down}}=1.661$ by a factor of $\sim2.7$, not a
percent-level match; the reciprocal convention ($1/|\Jones_{T(2,2)}|
=\vp$) does match to $2.6\%$, but this requires an unmotivated
choice of convention (see Remark~\ref{P19:rem:jones_ratio_coincidence}).
Whether the Jones-polynomial data of this paper has any bearing on
the observed ratio is open.
\end{proposition}
\status{draft}
\source{P19:prop:jones_modulus_ratio}
\depends{strong/confinement_topology.tex}

\begin{remark}[An unresolved numerical coincidence]
\label{P19:rem:jones_ratio_coincidence}
Using $1/|\Jones_{T(2,2)}(q_5)|=\vp$ in place of
$|\Jones_{T(2,2)}(q_5)|$ gives $\vp\approx1.618$, matching the
observed $1.661$ to $2.6\%$ --- the same match originally (and
incorrectly) claimed for this proposition. No physical justification
for using the reciprocal of the Jones modulus is established: the
solar mixing angle (Remark~\ref{P19:rem:solar_physical}) uses
$|\Jones_{T(2,2)}(q_5)|$ directly, not its reciprocal, for the same
quantity. Adopting the reciprocal convention only where it produces a
match is not a justified choice; it is recorded as an open,
unresolved coincidence, not a result.
\end{remark}
\status{draft}
\source{P19:rem:jones_ratio_coincidence}

\begin{remark}[Cross-ratio substructure and opposite monotonicity]
\label{P19:rem:cross_ratio}
The up-type and down-type generation gaps have opposite monotonicity
(up-type decreasing, down-type increasing); the crossed ratios
$\Delta n_{\mathrm{up}}^{(1)}/\Delta n_{\mathrm{down}}^{(2)}=1.678$ and
$\Delta n_{\mathrm{up}}^{(2)}/\Delta n_{\mathrm{down}}^{(1)}=1.639$
are both $\approx\vp$ (within $3.7\%$ and $1.3\%$), while the
same-index ratios are not: each up-type gap equals $\vp$ times the
complementary down-type gap.
\end{remark}
\status{draft}
\source{P19:rem:cross_ratio}

\begin{theorem}[The four quark mass-ratio gaps form a $\sqrt\vp$ geometric ladder]
\label{P19:thm:sqrt_phi_ladder}
Let $\Delta n_q^{(g)}=\log_\vp(m_{q,g+1}/m_{q,g})$. The four
intra-triplet generation gaps satisfy, to within $4\%$,
$\bigl(\Delta n_{\mathrm{down}}^{(1)},\Delta n_{\mathrm{down}}^{(2)},
\Delta n_{\mathrm{up}}^{(2)},\Delta n_{\mathrm{up}}^{(1)}\bigr)
=x\cdot(1,\sqrt\vp,\vp,\vp^{3/2})$, where
$x=\Delta n_{\mathrm{down}}^{(1)}=\log_\vp(m_s/m_d)$ is the base of
the ladder. The inversion ratio (larger gap/smaller gap) is
$\sqrt\vp=1.272$ for both triplets ($r_{\mathrm{up}}=1.2993$,
$r_{\mathrm{down}}=1.2691$, within $2.1\%$).
\end{theorem}
\status{draft}
\source{P19:thm:sqrt_phi_ladder}
\residual{up to 3.4\% (per-gap)}
\depends{strong/quark_masses.tex}

\begin{remark}[The strange quark as the algebraic anchor]
\label{P19:rem:strange_anchor}
The base of the ladder $x=\log_\vp(m_s/m_d)$ makes the strange quark
the unit of the ladder: $m_b/m_s=(m_s/m_d)^{\sqrt\vp}$ (0.7\%),
$m_t/m_c=(m_s/m_d)^{\vp}$ (6.1\%), $m_c/m_u=(m_s/m_d)^{\vp^{3/2}}$
(19\%). The prediction degrades from the base of the ladder upward,
consistent with Paper~XVI's always-positive, generation-growing
residual pattern.
\end{remark}
\status{draft}
\source{P19:rem:strange_anchor}

\begin{remark}[$m_s/m_d\approx n_e=20$: the spin-Casimir hypothesis]
\label{P19:rem:ms_md_integer}
The framework derives $m_s/m_d=\sqrt{k/2}\cdot e^{8\pi/9}=19.991$,
within $0.047\%$ of the integer $n_e=20$, the lepton CMB tower level
$n_e=2Q_{\mathrm{group}}^{(\ell)}=2\times10=20$
(Theorem~P12:thm:cmb\_identity). The factor $\sqrt{k/2}=
\sqrt{2C_2(\tfrac12)}$ at $k=3$ follows from the structural identity
$k=4C_2(\tfrac12)=3$, exact at $k=3$ and $j=\tfrac12$: it is the
spin-statistics correction from the quark being a spin-$\tfrac12$
fermion, correcting the bosonic E$_8$ T-phase formula
$m_s/m_d=e^{8\pi/9}=16.32$ (an $18\%$ undershoot on its own). A proof
that the fermionic character of quarks inserts exactly this factor,
closing the $0.047\%$ gap to the integer $20$, is Open
Problem~P19:OP6.
\end{remark}
\status{draft}
\source{P19:rem:ms_md_integer}
\residual{0.047\%}
\depends{strong/quark_masses.tex}

\begin{corollary}[Absolute quark mass predictions from the ladder]
\label{P19:cor:absolute_masses}
Combining Paper~XVI's derived $m_s=93.32\,\mathrm{MeV}$ (Proposition~%
P16:prop:strange\_match, $0.09\%$) with $m_s/m_d=19.991$
(Remark~\ref{P19:rem:ms_md_integer}) and
Theorem~\ref{P19:thm:sqrt_phi_ladder}, using only
framework-internal quantities (no PDG input):
$m_d=93.32/19.991=4.668\,\mathrm{MeV}$ (PDG $4.67$; $0.04\%$);
$m_b=93.32\times19.991^{\sqrt\vp}=4214\,\mathrm{MeV}$ (PDG $4183$;
$0.7\%$); $m_t/m_c=19.991^{\vp}=127$ (PDG $136$; $6\%$);
$m_c/m_u=19.991^{\vp^{3/2}}=476$ (PDG $589$; $19\%$).
\end{corollary}
\status{draft}
\source{P19:cor:absolute_masses}
\residual{0.04\%--19\% (four predictions)}
\depends{strong/quark_masses.tex; strong/topology_qh3.tex}

\begin{remark}[The bottom quark mass as a framework prediction]
\label{P19:rem:mb_prediction}
$m_b=4214\,\mathrm{MeV}$ ($0.7\%$) is entirely parameter-free (no
fitting, no free coefficients); using the integer $n_e=20$ in place
of $19.991$ gives $4216\,\mathrm{MeV}$ ($0.8\%$), a negligible change,
while using the pure E$_8$ prediction $e^{8\pi/9}=16.32$ (without the
spin-Casimir correction) gives $3256\,\mathrm{MeV}$ ($22\%$ off),
recovering Paper~XVI's systematic undershoot: the spin-Casimir factor
$\sqrt{k/2}$ is the key missing piece, not the accuracy of $m_s$.
\end{remark}
\status{draft}
\source{P19:rem:mb_prediction}

\begin{remark}[The physical origin of the $\sqrt\vp$ inversion ratio is open]
\label{P19:rem:sqrt_phi_meaning}
The inversion ratio $r=\sqrt\vp$ within each triplet
(Theorem~\ref{P19:thm:sqrt_phi_ladder}) is an empirical fit to PDG
mass data, independent of the Jones-polynomial computations of this
paper. The Jones polynomial modulus $|\Jones_{T(2,2)}(q_5)|=1/\vp$
does not equal the $\mathrm{SU}(2)_3$ quantum dimension $d_{1/2}=\vp$
--- they are reciprocals, not the same quantity --- so the
previously-proposed picture (quantum dimension as a ``probability
level,'' its square root as an ``amplitude level'' shared with the
Jones modulus) does not hold as stated. The physical origin of
$\sqrt\vp$ remains open.
\end{remark}
\status{draft}
\source{P19:rem:sqrt_phi_meaning}

\begin{remark}[The lepton solar angle and the Hopf-link Jones modulus]
\label{P19:rem:cross_sector_phi}
The leading-order solar mixing angle satisfies
$\tan\theta_{12}^{(0)}=1/\vp$ (Proposition~\ref{P19:prop:solar_angle_jones}),
via the $\mathrm{SU}(2)_3$ S-matrix ratio $S_{00}/S_{01}=1/\vp$. The
same modulus $1/\vp$ equals $|\Jones_{T(2,2)}(q_5)|$
(Theorem~\ref{P19:thm:channel_1}), a numerical coincidence recorded
but not established as a derivation. This is a single lepton-sector
coincidence; no corresponding quark-sector appearance of $\vp$ is
established (Proposition~\ref{P19:prop:jones_modulus_ratio} is open,
not proved). The PMNS correction scale $1/(k+2)=1/5$
(Paper~XVII, \S11.4) is the same $k+2=5$ entering the spin-Casimir
identity $k=4C_2(\tfrac12)$.
\end{remark}
\status{draft}
\source{P19:rem:cross_sector_phi}
\depends{quantum/neutrino.tex}

%────────────────────────────────────────────────────────────────────
% Two-triplet mass-scale ratio from z-tangent geometry
%────────────────────────────────────────────────────────────────────

\begin{definition}[Network $z$-momentum scale]
\label{P19:def:z_scale}
The characteristic $z$-momentum-squared scale of a site network is the
mean of the squared tangent $z$-component over the network,
$\langle|\tang_z|^2\rangle_{\mathrm{net}}=
\tfrac1{N_{\mathrm{net}}}\sum_{i\in\mathrm{net}}|\tang_z^{(i)}|^2$.
\end{definition}
\status{draft}
\source{P19:def:z_scale}

\begin{proposition}[Two-triplet scale ratio prediction]
\label{P19:prop:triplet_ratio}
Under Conjecture~\ref{P19:conj:six_quarks_reformulated}, the two-triplet
mass-scale ratio $M_{\mathrm{down}}/M_{\mathrm{up}}$ is predicted at
leading order by the down-type network $z$-momentum scale
(Definition~\ref{P19:def:z_scale}):
$M_{\mathrm{down}}/M_{\mathrm{up}}\sim
\langle|\tang_z|^2\rangle_{\mathrm{down}}=
\tfrac12(0.5249^2+0.321^2)=0.189$, since
$\langle|\tang_z|^2\rangle_{\mathrm{up}}=0$ at the crossings
(horizontal tangent, Theorem~P16:thm:crossing\_geometry).
\end{proposition}
\status{draft}
\source{P19:prop:triplet_ratio}
\depends{strong/topology_qh3.tex}

\begin{corollary}[Predicted vs.\ observed]
\label{P19:cor:triplet_agreement}
The prediction $0.189$ of Proposition~\ref{P19:prop:triplet_ratio}
matches the PDG geometric-mean observed value $0.157$ to $20\%$.
\end{corollary}
\status{draft}
\source{P19:cor:triplet_agreement}
\residual{20\%}

\begin{remark}[On the residual $20\%$ deviation]
\label{P19:rem:triplet_deviation}
The $20\%$ residual is consistent with the systematic undershooting
pattern Paper~XVI identifies for all quark mass predictions: the
leading-order geometric input is qualitatively correct but undershoots
by a factor attributable to a missing generation-dependent,
always-positive contribution.
\end{remark}
\status{draft}
\source{P19:rem:triplet_deviation}

\begin{remark}[QCD angular ordering applied to the twelve fraying sites]
\label{P19:rem:angular_ordering}
The twelve fraying sites cycle as $\cdots\to C\to A\to B\to A\to\cdots$
(A~=~crossing, B~=~midpoint, C~=~lobe); the total tangent rotation
across one segment ($C\to A\to B\to A$) is $177.44^\circ\approx180^\circ$,
the topological origin of back-to-back jet kinematics. Applying QCD
angular ordering (successive emissions at decreasing angle from the
threading site, $t=\pi/6$) to all eleven other sites gives a definite
emission sequence, with the first jet at $167.21^\circ$ from the
mirror crossing $t=11\pi/6$ (opening angle $12.79^\circ$, set by the
C-lobe curvature $\kappa_C=0.349$). Restricting to the three upper
A-crossings (the B-departure network), the two non-threading
crossings are \emph{exactly} $120^\circ$ from the threading tangent
(a $\Zthree$-symmetry consequence, equation~P19:eq:upper\_crossing\_angles),
confirming the cascade-model hierarchy of
Proposition~\ref{P18:prop:cascade_p18} but correcting it: the two
intermediate-energy crossings are exactly angle-degenerate, so the
$m_t$--$m_c$ splitting cannot come from the forward cascade alone and
requires the formation-event $\Zthree$-breaking azimuthal asymmetry.
The angular-ordering hierarchy (heaviest: threading-site region;
middle: the degenerate $120^\circ$ pair; lightest: the remaining lower
crossings) identifies top/charm/up with the corresponding sites;
deriving $m_t/m_c$ and $m_c/m_u$ quantitatively from the curvature
profile $\kappa(t)$ is an explicit target for the dynamic integrator
of Open Problem~O1.
\end{remark}
\status{draft}
\source{P19:rem:angular_ordering}
\depends{strong/topology_qh3.tex}

\begin{remark}[Angular ordering does not change the reactor angle or the quark mass predictions]
\label{P19:rem:angular_ordering_impact}
\textbf{Reactor angle:} unaffected directly. Paper~XVII's
$\sin^2\theta_{13}\sim3\varepsilon_M\sin^2\theta_M$ uses only fixed
geometric quantities ($\theta_M=31.66^\circ$, $\varepsilon_M\approx0.022$),
independent of emission ordering; the C-lobes are correctly excluded
from this formula (including them would overshoot the observed value
by a factor of $5$), since they emit into the baryon sector itself,
not across to the lepton sector. \textbf{Quark mass predictions:}
unaffected, for different reasons --- the two-triplet scale ratio
(Proposition~\ref{P19:prop:triplet_ratio}) uses fixed $z$-tangent
geometry; the $\sqrt\vp$ ladder (Theorem~P19:thm:sqrt\_phi\_ladder)
uses only accepted-as-exact PDG-empirical ratios, with no
ordering-dependent (or Jones-polynomial) input at all. The angular
ordering does sharpen the physical picture: the exact $120^\circ$
crossing degeneracy (Remark~\ref{P19:rem:angular_ordering}) shows the
top/charm mass splitting requires the formation-event azimuthal
asymmetry, not the ordering cascade alone.
\end{remark}
\status{draft}
\source{P19:rem:angular_ordering_impact}

%────────────────────────────────────────────────────────────────────
% Multi-component Jones polynomial framework
%────────────────────────────────────────────────────────────────────

\begin{lemma}[Disjoint-union factor at the WZW root]
\label{P19:lem:disjoint_factor}
At $q_5=e^{2\pi i/5}$, the Kauffman/Jones disjoint-union multiplier
$-(q_5^{1/2}+q_5^{-1/2})$ equals $-\vp$.
\end{lemma}
\status{draft}
\source{P19:lem:disjoint_factor}

\begin{proposition}[Two-component initial state factorisation]
\label{P19:prop:two_comp_initial}
For two disjoint $\QH=2$ Hopf links, $\Jones_{T(2,2)\sqcup T(2,2)}(q_5)
=-\vp\cdot\Jones_{T(2,2)}(q_5)^2$
(Lemma~\ref{P19:lem:disjoint_factor}), where
$\Jones_{T(2,2)}(q_5)=1+\zeta_5^3$, of modulus $1/\vp$
(Theorem~\ref{P19:thm:channel_1}). The product is irrational since
both factors are irrational.
\end{proposition}
\status{draft}
\source{P19:prop:two_comp_initial}
\depends{strong/confinement_topology.tex}

\begin{definition}[Effective Jones product]
\label{P19:def:effective_jones}
For a multi-component state $\bigsqcup_i L_i$ of $N$ disjoint
link-objects, the effective Jones product at $q_5$ is
$\Jones_{\mathrm{eff}}(\bigsqcup_i L_i)\equiv
(-\vp)^{N-1}\prod_{i=1}^N\Jones_{L_i}(q_5)\in\mathbb Z[\zeta_5]$,
where $\zeta_5=e^{2\pi i/5}$.
\end{definition}
\status{draft}
\source{P19:def:effective_jones}

\begin{theorem}[Multi-component reaction rule: unit equivalence]
\label{P19:thm:unit_reaction_rule}
A reaction $\bigsqcup_i L_i^{\mathrm{init}}\to
\bigsqcup_j L_j^{\mathrm{final}}$ is topologically permitted at the
WZW root of unity if and only if the effective Jones products
(Definition~\ref{P19:def:effective_jones}) of the initial and final
states differ by a unit in $\mathbb Z[\zeta_5]$ (a principal ideal
domain, class number~1 for $p=5$): units are
$\pm\zeta_5^k\vp^n$, $k\in\{0,\ldots,4\}$, $n\in\mathbb Z$. This is
strictly stronger than Paper~XVIII's rationality-class selection rule
(Theorem~P18:thm:jones\_rule): unit equivalence implies same
rationality class, but not conversely.
\end{theorem}
\status{draft}
\source{P19:thm:unit_reaction_rule}
\depends{strong/confinement_topology.tex}

\begin{remark}[Relation to Paper~XVIII's rationality-class rule]
\label{P19:rem:rationality_vs_unit}
Two elements in the same rationality class need not differ by a unit;
the unit-equivalence formulation of
Theorem~\ref{P19:thm:unit_reaction_rule} is the correct one for a
principal ideal domain.
\end{remark}
\status{draft}
\source{P19:rem:rationality_vs_unit}

\begin{corollary}[Recovery of Paper~XVIII's single-channel rules]
\label{P19:cor:single_channel_recovery}
In the single-component case $N_{\mathrm{init}}=N_{\mathrm{final}}=1$,
Theorem~\ref{P19:thm:unit_reaction_rule} reduces to
$\Jones_{L_{\mathrm{init}}}(q_5)/\Jones_{L_{\mathrm{final}}}(q_5)\in
\mathbb Z[\zeta_5]^\times$, recovering Paper~XVIII's single-channel
picture as the $N=1$ special case. Neither production channel
verified in this paper is itself an $N=1$ instance (both have
$N_{\mathrm{init}}\geq2$); no single-component $T(2,2)\to T(2,3)$
reaction is physically realised in this framework (the free-standing
trefoil's Jones value has norm $11$ in $\mathbb Z[\zeta_5]$, not a
unit), so this corollary records the algebraic reduction only.
\end{corollary}
\status{draft}
\source{P19:cor:single_channel_recovery}

\begin{theorem}[Channel~1 verification: two-tube reaction]
\label{P19:thm:channel_1}
The reaction $T(2,2)\sqcup T(2,2)\to T(2,3)\sqcup U$ satisfies the
unit-equivalence rule of Theorem~\ref{P19:thm:unit_reaction_rule},
with unit factor $u=\zeta_5+\zeta_5^4=1/\vp$ (real, norm $1$), using
$\Jones_{T(2,2)}(q_5)=1+\zeta_5^3$ (modulus $1/\vp$) and the
satellite value $\Jones_{\mathrm{sat}}(q_5)=q_5^3/\vp$.
\end{theorem}
\status{draft}
\source{P19:thm:channel_1}
\depends{strong/confinement_topology.tex}

\begin{theorem}[Channel~2 verification: satellite reaction]
\label{P19:thm:channel_2}
The reaction $T(2,2)\sqcup U\to T(2,3)$ satisfies the unit-equivalence
rule of Theorem~\ref{P19:thm:unit_reaction_rule}, with unit factor
$u=-\vp\cdot\zeta_5$ (norm $1$).
\end{theorem}
\status{draft}
\source{P19:thm:channel_2}
\depends{strong/confinement_topology.tex}

\begin{remark}[The unit factors have physical content]
\label{P19:rem:unit_meaning}
The unit factors $u_1=\zeta_5+\zeta_5^4=1/\vp$ (Channel~1) and
$u_2=-\vp\cdot\zeta_5$ (Channel~2) do \emph{not} differ by a pure
root of unity: their moduli are reciprocals, $|u_1|=1/\vp$,
$|u_2|=\vp$, so $u_2/u_1=-\vp^2\zeta_5$ --- an exact identity
differing from a pure phase by the real factor $\vp^2$. This
asymmetry traces to Channel~2 carrying one more power of the
disjoint-union factor $-\vp$ relative to its final state than
Channel~1 does relative to its own (the $(-\vp)^{N-1}$ term of
Definition~\ref{P19:def:effective_jones}, given the two channels'
different component-count changes). Whether this modulus asymmetry
has further physical meaning is open.
\end{remark}
\status{draft}
\source{P19:rem:unit_meaning}
\depends{strong/confinement_topology.tex}

%────────────────────────────────────────────────────────────────────
% Preliminary numerics beyond Q_H=3
%────────────────────────────────────────────────────────────────────

\begin{remark}[Preliminary numerics beyond $\QH=3$: no distinct equilibrium found]
\label{P19:rem:beyond_qh3_numerics}
Topology-protected constrained gradient flow (angle-clamp $2^\circ$
per grid point per step, following Paper~XVI's $\QH=3$ saddle-search
method) was run on the general $T(2,q_{\rm pol})$ construction at
three points spanning the unestablished $\QH=4$--$5$ region beyond
the trefoil: $q_{\rm pol}=3.764$ and $4.236$ (the golden-section
division points of $[3,5]$) and $q_{\rm pol}=4.0$ (a local energy
minimum of the unrelaxed construction, coinciding with the $\QH=4$
sector, which Paper~XVIII acknowledges as charge-conservation-permitted
but establishes no stability for, Remark~\ref{P18:rem:higher_charge_incoming}).
All three runs completed without triggering any dilution or
topology-loss halt condition, but none reached an energy plateau: all
three continued decreasing smoothly through the full step budget, at
comparable rates near the end of the run ($q_{\rm pol}=3.764$ and
$4.236$ reduced by a factor of $\sim12$; $q_{\rm pol}=4.0$ by a
factor of $\sim4.5$, retaining a substantially higher $J_4/K$). This
is consistent with all three points draining toward the same
topologically trivial vacuum at different rates, rather than any
settling near a distinct stable configuration --- the same
flat-bottomed, no-restoring-force picture already established for
$\QH=3$ itself (Remark~\ref{P18:rem:saddle_difficulty}) extending
into this unestablished region. This is a preliminary, unconverged
numerical observation, not a proof of instability at these
$q_{\rm pol}$ values; longer runs or a convergence-based stopping
criterion would be needed to settle the question.
\end{remark}
\status{draft}
\source{P19:rem:beyond_qh3_numerics}
\depends{strong/topology_qh3.tex}
